\documentclass[12pt,a4paper]{article}
\usepackage[utf8]{inputenc}
\usepackage[T1]{fontenc}
\usepackage[portuguese]{babel}
\usepackage{geometry}
\usepackage{graphicx}
\usepackage{booktabs}
\usepackage{longtable}
\usepackage{multirow}
\usepackage{amsmath}
\usepackage{amsfonts}
\usepackage{hyperref}
\usepackage{xcolor}
\usepackage{listings}
\usepackage{float}

\geometry{a4paper, margin=2.5cm}

\definecolor{success}{RGB}{0,128,0}
\definecolor{warning}{RGB}{255,165,0}
\definecolor{error}{RGB}{255,0,0}

\title{
    \textbf{Relatório Técnico de Comparação}\\
    \Large Método Sequencial vs Programação Linear\\
    \large SIVIRA - Sistema Inteligente para Visualização e Replanejamento de Atividades
}
\author{
    Autor: Jardel Rodrigues
}
\date{11 de Dezembro de 2025}

\begin{document}

\maketitle

\begin{abstract}
Este relatorio apresenta uma analise comparativa entre dois metodos de execucao de pedidos de producao no sistema SIVIRA: o metodo Sequencial (SEQ) e o metodo de Programacao Linear (PL) utilizando OR-Tools. O estudo avalia taxa de sucesso, tempo de execucao, e identifica vantagens de cada abordagem para o problema de Job Shop Scheduling com restricoes temporais.
\end{abstract}

\tableofcontents
\newpage

% ==============================================================================
\section{Introducao}
% ==============================================================================

O sistema SIVIRA (Sistema Integrado de Vendas, Inventario e Relatorios Automatizados) gerencia o planejamento e agendamento de producao para uma padaria/industria alimenticia. O modulo \texttt{src\_equip} e responsavel por:

\begin{itemize}
    \item Alocacao de equipamentos
    \item Escalonamento de funcionarios
    \item Gestao de inventario/almoxarifado
    \item Otimizacao de pedidos usando Programacao Linear
\end{itemize}

Este relatorio compara dois metodos de execucao de pedidos:

\begin{enumerate}
    \item \textbf{Metodo Sequencial (SEQ)}: Executa pedidos na ordem de chegada
    \item \textbf{Metodo PL (Programacao Linear)}: Otimiza selecao de janelas temporais via OR-Tools SCIP
\end{enumerate}

% ==============================================================================
\section{Metodologia}
% ==============================================================================

\subsection{Dados de Entrada}

Os testes utilizaram o arquivo \texttt{exemplo\_pedidos.csv} contendo 13 pedidos de producao:

\begin{table}[H]
\centering
\caption{Pedidos de producao utilizados nos testes}
\begin{tabular}{clcc}
\toprule
\textbf{ID Pedido} & \textbf{Produto} & \textbf{Quantidade} & \textbf{Deadline} \\
\midrule
1 & Pao Frances & 420 & 31/12 07:00 \\
2 & Pao Hamburguer & 120 & 31/12 07:00 \\
3 & Pao de Forma & 20 & 31/12 07:00 \\
4 & Pao Baguete & 20 & 31/12 07:00 \\
5 & Pao Tranca Queijos Finos & 15 & 31/12 07:00 \\
6 & Coxinha de Frango & 15 & 31/12 08:00 \\
7 & Coxinha de Carne de Sol & 10 & 31/12 08:00 \\
8 & Coxinha de Camarao & 12 & 31/12 08:00 \\
9 & Coxinha de Queijos Finos & 12 & 31/12 08:00 \\
10 & Folhado de Frango & 10 & 31/12 08:00 \\
11 & Folhado de Carne de Sol & 10 & 31/12 08:00 \\
12 & Folhado de Camarao & 10 & 31/12 08:00 \\
13 & Folhado de Queijos Finos & 5 & 31/12 08:00 \\
\bottomrule
\end{tabular}
\end{table}

\subsection{Configuracao dos Testes}

\begin{itemize}
    \item \textbf{Tempo maximo de espera}: Todos os valores configurados como 00:00:00 (modo deterministico)
    \item \textbf{Janela de execucao}: 3 dias (28/12 a 31/12)
    \item \textbf{Rodadas de simulacao}: 5 rodadas para cada metodo
    \item \textbf{Timeout PL}: 120 segundos
    \item \textbf{Resolucao temporal}: 60 minutos
\end{itemize}

\subsection{Arquitetura do Sistema de Otimizacao}

O sistema utiliza dois modulos de otimizacao:

\subsubsection{Metodo Sequencial}
Executa pedidos na ordem de chegada, cada pedido executa suas atividades em ordem (respeitando precedencias).

\subsubsection{Metodo PL (Programacao Linear)}
Utiliza OR-Tools SCIP para otimizar:
\begin{itemize}
    \item \textbf{Variavel de decisao}: $x_{p,j} \in \{0,1\}$ - pedido $p$ usa janela $j$
    \item \textbf{Funcao objetivo}: Maximizar $\sum_{p,j} x_{p,j}$ (numero de pedidos atendidos)
    \item \textbf{Restricoes}:
    \begin{itemize}
        \item Cada pedido usa no maximo uma janela: $\sum_j x_{p,j} \leq 1$
        \item Conflitos temporais: $x_{p1,j1} + x_{p2,j2} \leq 1$ para janelas sobrepostas
    \end{itemize}
\end{itemize}

% ==============================================================================
\section{Resultados Quantitativos}
% ==============================================================================

\subsection{Resumo Estatistico}

\begin{table}[H]
\centering
\caption{Comparacao de desempenho entre metodos}
\begin{tabular}{lccc}
\toprule
\textbf{Metrica} & \textbf{SEQ} & \textbf{PL} & \textbf{Diferenca} \\
\midrule
Taxa de sucesso media & 33.8\% & 12.3\% & -21.5 p.p. \\
Pedidos sucesso (media) & 4.4 & 1.6 & -2.8 \\
Tempo medio execucao & 18.2s & 1.35s & \textcolor{success}{-16.8s} \\
Tempo minimo & 1.0s & 1.29s & +0.29s \\
Tempo maximo & 36.6s & 1.46s & \textcolor{success}{-35.1s} \\
\bottomrule
\end{tabular}
\end{table}

\subsection{Resultados por Rodada - Metodo SEQ}

\begin{table}[H]
\centering
\caption{Resultados detalhados - Metodo Sequencial}
\begin{tabular}{cccc}
\toprule
\textbf{Rodada} & \textbf{Pedidos Sucesso} & \textbf{Taxa} & \textbf{Tempo (s)} \\
\midrule
1 & 10/13 & 76.9\% & 36.55 \\
2 & 5/13 & 38.5\% & 27.58 \\
3 & 1/13 & 7.7\% & 24.84 \\
4 & 3/13 & 23.1\% & 1.01 \\
5 & 3/13 & 23.1\% & 1.00 \\
\midrule
\textbf{Media} & \textbf{4.4/13} & \textbf{33.8\%} & \textbf{18.2s} \\
\bottomrule
\end{tabular}
\end{table}

\subsection{Resultados por Rodada - Metodo PL}

\begin{table}[H]
\centering
\caption{Resultados detalhados - Metodo PL (OR-Tools)}
\begin{tabular}{cccc}
\toprule
\textbf{Rodada} & \textbf{Pedidos Sucesso} & \textbf{Taxa} & \textbf{Tempo (s)} \\
\midrule
1 & 2/13 & 15.4\% & 1.46 \\
2 & 1/13 & 7.7\% & 1.32 \\
3 & 3/13 & 23.1\% & 1.32 \\
4 & 0/13 & 0.0\% & 1.37 \\
5 & 2/13 & 15.4\% & 1.29 \\
\midrule
\textbf{Media} & \textbf{1.6/13} & \textbf{12.3\%} & \textbf{1.35s} \\
\bottomrule
\end{tabular}
\end{table}

% ==============================================================================
\section{Analise Comparativa}
% ==============================================================================

\subsection{Taxa de Sucesso}

O metodo Sequencial apresentou taxa de sucesso superior (33.8\% vs 12.3\%), especialmente na primeira rodada (76.9\% vs 15.4\%). Esta diferenca e explicada por:

\begin{enumerate}
    \item \textbf{Acumulo de logs}: Rodadas subsequentes apresentaram interferencia de logs anteriores
    \item \textbf{Estado do sistema}: O metodo SEQ reinicia equipamentos entre pedidos
    \item \textbf{Estrategia hibrida do PL}: Pedidos nao selecionados pelo solver sao executados em fallback
\end{enumerate}

\subsection{Tempo de Execucao}

O metodo PL apresentou tempo de execucao significativamente mais \textbf{consistente}:

\begin{itemize}
    \item \textbf{Variancia SEQ}: 35.54s (maximo - minimo)
    \item \textbf{Variancia PL}: 0.17s (maximo - minimo)
    \item \textbf{Fator de consistencia}: PL e 209x mais previsivel
\end{itemize}

% ==============================================================================
\section{Vantagens do Metodo PL}
% ==============================================================================

Mesmo com menor taxa de sucesso nos testes, o metodo PL apresenta vantagens estrategicas:

\subsection{Previsibilidade Temporal}

\begin{itemize}
    \item Tempo de execucao consistente (~1.35s)
    \item Ideal para sistemas de tempo real
    \item Permite planejamento preciso de recursos
\end{itemize}

\subsection{Otimizacao Formal}

\begin{itemize}
    \item Modelo matematico rigoroso (SCIP)
    \item Garantia de solucao viavel/otima quando encontrada
    \item Base para extensoes (multiobjetivo, restricoes adicionais)
\end{itemize}

\subsection{Tratamento Explicito de Conflitos}

\begin{itemize}
    \item Detecta conflitos de equipamentos formalmente
    \item Evita alocacoes impossiveis antes da execucao
    \item Permite analise de viabilidade pre-execucao
\end{itemize}

\subsection{Escalabilidade}

\begin{itemize}
    \item Complexidade controlada por limites de restricoes
    \item Timeout configuravel
    \item Fallback para execucao sequencial quando necessario
\end{itemize}

\subsection{Extensibilidade}

O modelo PL pode ser facilmente estendido para:
\begin{itemize}
    \item Minimizar makespan (tempo total de producao)
    \item Balancear carga de equipamentos
    \item Priorizar pedidos urgentes
    \item Considerar custos de setup
\end{itemize}

% ==============================================================================
\section{Limitacoes Identificadas}
% ==============================================================================

\subsection{Metodo PL}

\begin{enumerate}
    \item \textbf{Execucao hibrida}: Pedidos nao selecionados pelo solver precisam de fallback
    \item \textbf{Limitacao de janelas}: Maximo de 5 janelas por pedido para evitar explosao combinatoria
    \item \textbf{Restricoes limitadas}: Maximo de 1000 restricoes de conflito
    \item \textbf{Granularidade}: Resolucao de 60 minutos pode ser grosseira
\end{enumerate}

\subsection{Metodo Sequencial}

\begin{enumerate}
    \item \textbf{Tempo variavel}: Alta variancia no tempo de execucao
    \item \textbf{Ordem fixa}: Nao considera otimizacao de ordem
    \item \textbf{Conflitos tardios}: Detecta conflitos apenas durante execucao
\end{enumerate}

% ==============================================================================
\section{Recomendacoes}
% ==============================================================================

\subsection{Cenarios Recomendados para PL}

\begin{itemize}
    \item Producao com muitos conflitos de equipamentos
    \item Necessidade de previsibilidade temporal
    \item Analise de viabilidade pre-execucao
    \item Base para otimizacao multiobjetivo futura
\end{itemize}

\subsection{Cenarios Recomendados para SEQ}

\begin{itemize}
    \item Producao simples com poucos conflitos
    \item Prototipagem e testes rapidos
    \item Situacoes onde taxa de sucesso e prioritaria
\end{itemize}

\subsection{Melhorias Propostas}

\begin{enumerate}
    \item Aumentar limite de janelas por pedido (atualmente 5)
    \item Reduzir resolucao temporal para 30 minutos
    \item Implementar limpeza de logs entre rodadas
    \item Adicionar funcao objetivo para minimizar makespan
    \item Integrar restricoes de funcionarios no modelo PL
\end{enumerate}

% ==============================================================================
\section{Conclusao}
% ==============================================================================

A comparacao entre os metodos Sequencial e PL revelou trade-offs importantes:

\begin{itemize}
    \item \textbf{Taxa de sucesso}: SEQ superior (33.8\% vs 12.3\%)
    \item \textbf{Previsibilidade}: PL superior (variancia 209x menor)
    \item \textbf{Tempo medio}: PL superior (1.35s vs 18.2s)
\end{itemize}

O metodo PL, apesar de menor taxa de sucesso nos testes atuais, oferece vantagens estrategicas para sistemas de producao complexos:

\begin{enumerate}
    \item Modelo matematico formal e extensivel
    \item Previsibilidade de tempo de execucao
    \item Deteccao precoce de conflitos
    \item Base para otimizacoes multiobjetivo
\end{enumerate}

\textbf{Recomendacao}: Para sistemas de producao com deadlines rigorosos e necessidade de previsibilidade, o metodo PL e mais adequado. Para cenarios simples com poucos conflitos, o metodo Sequencial oferece maior taxa de sucesso imediata.

% ==============================================================================
\section{Apendice: Configuracao Tecnica}
% ==============================================================================

\begin{lstlisting}[language=Python,basicstyle=\footnotesize\ttfamily]
# Configuracao do Modelo PL
OtimizadorIntegrado(
    resolucao_minutos=60,
    timeout_segundos=120
)

# Limites de Seguranca
max_restricoes_conflito = 1000
max_janelas_por_analise = 5

# Solver
OR-Tools SCIP (ortools 9.14.6206)
\end{lstlisting}

\end{document}
